\documentclass[14pt,a4paper]{extarticle}

\usepackage[utf8]{inputenc}
\usepackage[english]{babel}
\usepackage{setspace}
\usepackage{geometry}

\geometry{
    left=25mm,
    right=15mm,
    top=20mm,
    bottom=20mm
}

\setstretch{1.5}

\begin{document}

\section*{Abstract}

\textbf{Aleksieienko~A.\,D.} Optimization of Genetic Algorithm Computations through CUDA Parallel Processing. - Course paper. National University of ``Kyiv-Mohyla Academy'', Faculty of Informatics, Department of Computer Science. - Kyiv, 2026.

\bigskip
This paper investigates the acceleration of genetic algorithms (GAs) through CUDA-based GPU parallelization using Python-native Numba JIT-compiled kernels. Custom kernels for fitness evaluation, selection, crossover, and mutation were benchmarked against CPU implementations on XOR perceptron training and the Traveling Salesman Problem.

GPU implementations outperform CPU baselines for populations exceeding 100 individuals, achieving speedups of 1.15$\times$ to 14.56$\times$ on an NVIDIA RTX~2050. All differences are statistically significant ($p < 0.05$, Cohen's $d > 0.8$). The dominant bottleneck is host-device data transfer, indicating further gains are achievable with fully GPU-resident evolutionary loops. The codebase was refactored into a publicly available Python library for GPU-accelerated GA development.

\bigskip
\textbf{Keywords:} genetic algorithms, CUDA, GPU parallelization, Numba, performance optimization, traveling salesman problem, neural network training, Python, parallel computing.

\end{document}
